\section{Complete characteristic series and topology}
\label{final:series}

The complete characteristic series below recover the intrinsic
decomposition of $\Sigma$. The coordinate in every formal statement
is marked. Unless
otherwise specified, coefficients lie in a field $K$ of characteristic
zero.  The same identities hold in a commutative $\mathbb Q$-algebra
when every scalar used as a denominator is a unit.  For a formal series,
$[u^n]$ denotes coefficient extraction and $\operatorname{Res}_u$ is
the coefficient of $u^{-1}$ in a Laurent differential.  A quotient
such as $u/(1-e^{-u})$ means the inverse of the unit
$(1-e^{-u})/u$; no inverse of the nonunit $u$ is being asserted.

\begin{theorem}[The normalized Todd tower and all its twists]
\label{final:F1}
There is a unique series $Q(u)=\sum_{j\geq0}q_ju^j$ satisfying
\[
 q_0=1,\qquad [u^n]Q(u)^{n+1}=1\quad(n\geq1).
\]
It is
\[
 Q_T(u)=\frac{u}{1-e^{-u}}.
\]
Equivalently, one may impose the coefficient equations for every
$n\geq0$, since their degree-zero equation is exactly $q_0=1$.
For this series, for every $n\geq0$,
\begin{equation}\label{final:todd-twists}
 [u^n]e^{ku}Q_T(u)^{n+1}
   =\binom{n+k}{n}
   =\frac{(k+1)(k+2)\cdots(k+n)}{n!}.
\end{equation}
This is a polynomial identity in $k$ over $\mathbb Q$, including the
empty-product convention for $n=0$, and hence holds at every $k\in K$.
In particular it holds for all positive and negative integer twists.
Every individual positive-degree equation is indispensable within the
displayed tower with $q_0=1$.
\end{theorem}

\begin{proof}
For $n\geq1$, expanding the product of $n+1$ copies of $Q$ gives
\[
 [u^n]Q^{n+1}=(n+1)q_n+P_n(q_1,\ldots,q_{n-1}).
\]
The only occurrences of $q_n$ choose degree $n$ from one factor and
degree zero from the others.  As $n+1$ is invertible, this equation
determines $q_n$ from the preceding coefficients.  This proves
uniqueness, provided existence is established.

Put $z=1-e^{-u}$, whose compositional inverse is $u=-\log(1-z)$.
The formal change-of-coordinate rule for residues gives
\begin{align*}
 [u^n]e^{ku}Q_T^{n+1}
 &=\operatorname{Res}_u\frac{e^{ku}}{(1-e^{-u})^{n+1}}\,du\\
 &=\operatorname{Res}_z z^{-n-1}(1-z)^{-k-1}\,dz
   =[z^n](1-z)^{-k-1}
   =\binom{n+k}{n}.
\end{align*}
For completeness, residue substitution by a coordinate
$u=\phi(z)=z+O(z^2)$ follows by linearity from Laurent monomials.
If $j\ne-1$, $\phi^j\phi'$ is the derivative of
$\phi^{j+1}/(j+1)$ and has zero residue; if $j=-1$, write
$\phi=zv$, $v(0)=1$, to obtain
$\phi'/\phi=z^{-1}+v'/v$, with residue one.  Only finitely many
negative Laurent powers are involved, and nonnegative powers cannot
contribute a residue.  The final coefficient identity follows from
the binomial series, or by solving
$(1-z)F'=(k+1)F$, $F(0)=1$, coefficient by coefficient.
Taking $k=0$ proves existence and all tower equations.

If the equation at a positive degree $N$ is removed, retain the
uniquely determined coefficients below $N$, choose any $q_N$ different
from its Todd value, and apply the same recursion at every $n>N$.
The resulting series satisfies every retained equation and differs
from $Q_T$.  A finite initial tower leaves every later coefficient
unconstrained.  Finally, if $q_0$ is only required to be a unit and the
degree-zero equation is absent, the coefficient of $q_n$ becomes
$(n+1)q_0^n$.  Thus each chosen unit $q_0$ gives a recursively determined
tower; normalization cannot be silently dropped.
\end{proof}

For comparison, arbitrary powers have the separate finite coefficient
formula
\begin{equation}\label{final:todd-arbitrary-power}
 [u^n]Q_T(u)^k=
 \sum_{\substack{m_1,\ldots,m_n\geq0\\\sum j m_j=n}}
 (k)_{\underline M}\prod_{j=1}^n\frac{q_j^{m_j}}{m_j!},
 \qquad M=\sum_{j=1}^n m_j,
\end{equation}
where $(k)_{\underline M}=k(k-1)\cdots(k-M+1)$ and the
empty product is one.  Indeed, apply the formal binomial expansion
to $Q_T=1+(Q_T-1)$ and then the multinomial theorem.  This formula
does not replace the exponent $n+1$ and exponential twist in
\eqref{final:todd-twists}.

\begin{theorem}[Reversible characteristic-series formulas]
\label{final:F2}
Define the normalized series
\[
 Q_A(u)=\frac{u/2}{\sinh(u/2)},\qquad
 Q_L(u)=\frac{u}{\tanh u},\qquad
 Q_y(u)=Q_T((1+y)u)-yu.
\]
Their constant terms are one.  With $y$ fixed and marked and
$1+y\ne0$, each of $Q_T,Q_A,Q_L,Q_y$, and each of their multiplicative
inverses, determines the others in the marked formal coordinate.
Explicit reconstruction formulas are
\begin{gather}
 Q_T=e^{u/2}Q_A,\qquad Q_L(u)=Q_T(2u)-u,
 \qquad Q_T(u)=Q_L(u/2)+u/2,\label{final:TAL}\\
 e^u=\frac{Q_T(u)}{Q_T(u)-u},\qquad
 e^{2u}=\frac{Q_L(u)+u}{Q_L(u)-u},\label{final:exponential-recovery}\\
 Q_T(v)=Q_y\left(\frac{v}{1+y}\right)+\frac{yv}{1+y}.
 \label{final:chi-inverse}
\end{gather}
For a reconstruction from $Q_A$ alone, set
\[
 b=\frac{u}{Q_A(u)},\qquad
 w=\frac{b+\sqrt{b^2+4}}{2},\qquad Q_T=wQ_A,
\]
where the formal square root has constant term $2$.
The specialization $y=-1$ is exactly $Q_{-1}(u)=1+u$ and has no
inverse of the form claimed.  Over a general $\mathbb Q$-algebra the
hypothesis is that $1+y$ is a unit, and over $\mathbb Q[y]$ the
inverse formula takes place after localization at $1+y$.
\end{theorem}

\begin{proof}
Substitution of $\sinh v=(e^v-e^{-v})/2$ and
$\tanh v=(e^{2v}-1)/(e^{2v}+1)$ proves
\eqref{final:TAL}--\eqref{final:exponential-recovery}.  All denominators
in the displayed ratios are units.  The square-root recursion is
unique because its leading coefficient is $2$, so its successive
unknown coefficients have invertible coefficient $4$.  In the present
case $b=2\sinh(u/2)$ and this root is $2\cosh(u/2)$; thus $w=e^{u/2}$.
Equivalently $w(0)=1$ and $w-w^{-1}=b$, whose unit solution is unique.
The affine rescaling defining $Q_y$ gives \eqref{final:chi-inverse}
directly, and $y=-1$ gives $1+u$.  Multiplicative inversion on units
is an involution.  These calculations establish the stated reversals
on the exact series; applying the square-root algorithm to an
arbitrary even unit does not certify that its input equals $Q_A$.
\end{proof}

\begin{proposition}[Formal, analytic, and global recovery]
\label{final:F3}
The preceding reversals recover the formal exponential and its inverse
formal logarithm, hence the germ of $H(1+z)=\log(1+z)-z$ at $z=0$.
After specializing coefficients and the marked parameter $y$ to
the reals, they also construct the displayed standard real analytic functions
on their full real domains.  They identify an independently supplied
global function from its germ only in a declared uniqueness category,
for example real analytic functions on the same connected domain.
No analogous identification of arbitrary smooth global extensions is
valid.
\end{proposition}

\begin{proof}
Equation \eqref{final:exponential-recovery} recovers $e^u$, and the
series $e^u-1$ has linear coefficient one, so it has a unique
compositional inverse, $\log(1+z)$.  The real quotient definitions
extend through $u=0$ with value one.  Their other denominators do not
vanish on the real line; in the $Q_A$ reconstruction the chosen root
is $2\cosh(u/2)>0$.  Consequently the explicit formulas define the
claimed real analytic extensions.  The analytic identity theorem on
a connected interval proves uniqueness within that category.

The nearest nonremovable complex poles are $\pm2\pi i$ for $Q_T,Q_A$
and $\pm\pi i$ for $Q_L$.  Their Taylor discs therefore do not themselves
cover the real line.  For the smooth counterexample, choose a nonzero
smooth bump $\chi$ supported in a compact subinterval of $(0,1)$ away
from $1$.  Then $I+\varepsilon\chi$ has the same germ at $1$ as $I$
and differs globally for $\varepsilon\ne0$.  Since $I''=t^{-2}$ has a
positive minimum on the support, a sufficiently small
$|\varepsilon|$ preserves strict convexity.  A formal identity alone
cannot exclude this example.
\end{proof}

\begin{theorem}[The universal Thom correction, conditional on topology]
\label{final:F4}
Suppose complex topological $K$-theory and rational ordinary cohomology
are supplied with their Thom isomorphisms, compatible complex
orientations, natural multiplicative Thom classes, the Chern character
normalized by $\operatorname{ch}(L)=e^{c_1(L)}$, and an injective
splitting principle.  For a complex vector bundle $V\to B$ on a finite
CW complex, with projection $\pi$, the uniquely defined correction in
\begin{equation}\label{final:Thom-comparison}
 \operatorname{ch}(U_K(V))=\pi^*C(V)\,U_H(V)
\end{equation}
is
\[
 C(V)=\operatorname{Td}(V)^{-1}
      =\prod_i Q_T(u_i)^{-1},
\]
where the $u_i$ are formal Chern roots.  Its complete universal line
series determines $Q_T$ and conversely.  With compatible collapse and
Gysin constructions supplied as well, this yields
\begin{equation}\label{final:RR}
 \operatorname{ch}(f_!^K a)
   =f_*^H\bigl(\operatorname{ch}(a)\operatorname{Td}(T_f)\bigr)
\end{equation}
for $a\in K^0$ and a proper smooth map between manifolds with the
specified complex $K$-orientation and virtual relative tangent bundle
$T_f$, using compact supports when required.  On universal bases
the characteristic-class statement is taken in completed rational
cohomology.
\end{theorem}

\begin{proof}
The cohomological Thom isomorphism makes multiplication by $U_H(V)$
an isomorphism from the cohomology of the base to the appropriate
relative group, so it gives existence and uniqueness of $C(V)$ in
\eqref{final:Thom-comparison}.  For a line bundle put $u=c_1(L)$.
The zero-section pullbacks of the two Thom classes are respectively
$1-[L^*]$ and $u$.  Thus
\[
 uC(u)=1-e^{-u}.
\]
Naturality makes these identities compatible on the finite stages
of the universal line bundle.  Their inverse limit is
$\mathbb Q[[u]]$, in which multiplication by $u$ is injective.
Hence $C(u)=(1-e^{-u})/u=Q_T(u)^{-1}$.  Multiplicativity of Thom
classes and the Chern character gives the product formula after
splitting, and injectivity of the supplied splitting map descends
it to $B$.  For an embedding, applying the collapse map gives the
cohomological pushforward with the correction
$\operatorname{Td}(\nu)^{-1}$ for its complex normal bundle $\nu$.
Since $T_f=-\nu$, this is \eqref{final:RR}; the supplied stable
factorization and composition rules give the same formula for the
stated proper maps.  Finally, the universal line correction is a
unit, so inversion recovers $Q_T$.
\end{proof}

The scalar Euler-coordinate identity used here is itself elementary:
for $x\star y=x+y+xy$ and $z(x)=x/(1+x)$,
\[
 z(x\star y)=z(x)+z(y)-z(x)z(y).
\]
Indeed $1-z(x)=1/(1+x)$ is multiplicative under $\star$.
Interpreting $x$ as $[L]-1$ and $z$ as $1-[L^*]$ is a map into a
supplied $K$-theory ring.  The scalar identity does not construct
bundles, Thom isomorphisms, or the underlying spaces.  On a finite
base $u$ may be nilpotent and a zero divisor; cancellation in that
base would be invalid.  If $u^{N+1}=0$, adding $u^{N+1}R(u)$ to a
characteristic series does not change its evaluation there.  A single
finite-base Thom comparison or numerical index consequently cannot
determine the complete universal series.  Interpretations of the
$L$, $\widehat A$, and $\chi_y$ series retain their respective
oriented, spin, or complex/holomorphic foundations; the algebraic
formulas do not assert those foundations or their index theorems.

\begin{proposition}[Rational characteristic data lose integral data]
\label{final:F5}
On $\mathbb{RP}^2$, the complexification $L$ of the tautological real
line is nontrivial and $[L]\ne[1]$ in complex $K^0$, although
$\operatorname{ch}(L)=1$ in rational cohomology.  Thus even with the
base fixed, rational characteristic data and the complete scalar
series do not reconstruct integral bundle or $K$-theory data.
\end{proposition}

\begin{proof}
The usual cellular calculation gives
$H^2(\mathbb{RP}^2;\mathbb Z)=\mathbb Z/2$ and no positive-degree
rational cohomology.  To verify nontriviality directly, a nowhere-zero
section of $L$ would pull back to an odd map
$g:S^2\to\mathbb C\setminus\{0\}$.  Normalize it to an odd map
$S^2\to S^1$.  Since $S^2$ is simply connected, write
$g(x)=e^{i\theta(x)}$ for a continuous real function $\theta$.
Then $\theta(-x)-\theta(x)$ is a continuous odd integral multiple
of $\pi$, hence constant by connectedness.  Replacing $x$ by $-x$
changes its sign, which is impossible for a nonzero odd multiple
of $\pi$.  Thus $L$ is nontrivial.  Equivalently, under the usual
classification of complex line bundles, its first Chern class is
the nonzero element of the displayed group.

If $[L]=[1]$, the Grothendieck-group relation supplies a bundle $E$
with $L\oplus E\cong 1\oplus E$.  Taking determinant lines and
cancelling $\det E$ would make $L$ trivial, a contradiction.
On the other hand, every positive-degree rational Chern class
vanishes, so $\operatorname{ch}(L)=1$.  The transition functions of
$L$ are signs, so $L\otimes L\cong1$, as expected from the torsion
line class.  These arguments use the stated elementary topological
foundations and do not follow from formal scalar arithmetic alone.
\end{proof}

\section{The matrix lift and its Gaussian statistics}
\label{final:matrix}

Write $\mathcal P_n$ for the cone of real symmetric positive-definite
$n\times n$ matrices, $E_n$ for its identity, and use the trace pairing
$\langle U,V\rangle=\operatorname{tr}(UV)$ on symmetric matrices.
All matrix logarithms, square roots, and powers below are the positive
spectral ones.  The finite-dimensional real spectral theorem is
retained throughout this section.

\begin{theorem}[The unique scalar-block spectral lift]
\label{final:M1}
Suppose a family $F_n:\mathcal P_n\to\mathbb R$, $n\geq1$, satisfies
\[
 F_1([t])=I(t),\qquad F_n(QXQ^{\mathsf T})=F_n(X)
\]
for orthogonal $Q$, and the scalar-block recursion
\[
 F_{n+1}(X\oplus[t])=F_n(X)+I(t).
\]
Then, without any regularity assumption,
\begin{equation}\label{final:matrix-potential}
 F_n(X)=\Phi_n(X):=\sum_{i=1}^n I(\lambda_i(X))
          =\operatorname{tr}X-\log\det X-n.
\end{equation}
This family satisfies full block additivity.  Neither orthogonal
invariance nor scalar-block recursion can be removed while keeping
the scalar seed, even among nonnegative candidates.
\end{theorem}

\begin{proof}
Orthogonally diagonalize $X$ and apply scalar-block recursion until
all blocks have size one.  This forces the sum in
\eqref{final:matrix-potential}.  Trace, determinant, and their behavior
on block diagonal matrices verify every axiom for $\Phi_n$.

If orthogonal invariance is omitted,
$F_n(X)=\sum_i I(X_{ii})$ retains the seed and full block additivity.
Rotating $\operatorname{diag}(2,1)$ through angle $\pi/4$ gives diagonal
entries $3/2,3/2$, and
$I(2)-2I(3/2)=\log(9/8)\ne0$, so it is a different lift.
If scalar-block recursion is omitted, for $\varepsilon>0$ put
\[
 G_n(X)=\Phi_n(X)+\varepsilon
   \sum_{i<j}(\log\lambda_i(X)-\log\lambda_j(X))^2.
\]
This has the seed, orthogonal invariance, and nonnegativity, but on
$[a]\oplus[b]$, $a\ne b$, its extra term is nonzero whereas both
rank-one extra terms vanish.  If the seed is omitted, the zero family
already retains the two structural rules.  These are omission tests
in the stated family, rather than a claim of an absolute weakest
possible axiomatization.
\end{proof}

\begin{theorem}[Derivatives, divergence, duality, and determinant bounds]
\label{final:M2}
In the indicated affine coordinates the potential $\Phi_n$ satisfies
\begin{align}
 \nabla\Phi_n(X)&=E_n-X^{-1},\label{final:matrix-gradient}\\
 g_X(U,V):=D^2\Phi_n(X)[U,V]
   &=\operatorname{tr}(X^{-1}UX^{-1}V),\label{final:matrix-hessian}\\
 D(X,Y):=\Phi_n(X)-\Phi_n(Y)
   -\langle\nabla\Phi_n(Y),X-Y\rangle
   &=\operatorname{tr}(Y^{-1}X)-\log\det(Y^{-1}X)-n.
   \label{final:matrix-divergence}
\end{align}
The Hessian is positive definite and $D(X,Y)\geq0$, with equality
exactly when $X=Y$.  For every invertible real $C$, simultaneous
congruence preserves both $D$ and $g$.  Inversion is an isometry for
$g$ and reverses the directed divergence:
\[
 D(X^{-1},Y^{-1})=D(Y,X).
\]
Extending $\Phi_n$ by $+\infty$ outside $\mathcal P_n$ in the vector
space of symmetric matrices, its conjugate is
\begin{equation}\label{final:matrix-conjugate}
 \Phi_n^*(\Theta)=
 \begin{cases}
  -\log\det(E_n-\Theta),&E_n-\Theta\in\mathcal P_n,\\
  +\infty,&\text{otherwise}.
 \end{cases}
\end{equation}
The Legendre coordinate is $\Theta=E_n-X^{-1}$, with inverse
$X=(E_n-\Theta)^{-1}$.  Moreover
\[
 \det X\leq e^{\operatorname{tr}X-n},\qquad
 \det X\leq\left(\frac{\operatorname{tr}X}{n}\right)^n,
\]
with equality respectively at $E_n$ and at positive scalar matrices;
$\log\det$ is strictly concave.
\end{theorem}

\begin{proof}
Jacobi's formula gives $D\log\det X[U]=\operatorname{tr}(X^{-1}U)$.
Differentiating $XX^{-1}=E_n$ gives
$D(X^{-1})[U]=-X^{-1}UX^{-1}$ and proves the first two formulas.
Their positivity is the precise identity
\[
 g_X(U,U)=\operatorname{tr}
    \bigl((X^{-1/2}UX^{-1/2})^2\bigr)
          =\|X^{-1/2}UX^{-1/2}\|_F^2.
\]
It is the trace of a square, and is strictly positive for $U\ne0$.
Substitution yields \eqref{final:matrix-divergence}.  With
$A=Y^{-1/2}XY^{-1/2}$, trace and determinant show
$D(X,Y)=\sum_i I(\lambda_i(A))$.  Scalar nonnegativity and its equality
case prove the assertion, since $A=E_n$ precisely when $X=Y$.
Cyclic trace and the inverse-congruence formula prove congruence
invariance.  Substituting inverses in $D$ gives its reversal; inserting
the differential $-X^{-1}UX^{-1}$ in $g_{X^{-1}}$ gives $g_X$.

If $B=E_n-\Theta$ is positive definite, maximizing the conjugate is
equivalent to maximizing $-\operatorname{tr}(BX)+\log\det X+n$.
Substitute $A=B^{1/2}XB^{1/2}$ to obtain
$-\Phi_n(A)-\log\det B$.  Its unique maximum occurs at $A=E_n$,
which proves the finite formula and its optimizer.  If $B$ has an
eigenvalue $b\leq0$, take $X$ diagonal in an eigenbasis of $B$, with
the matching eigenvalue equal to $r\to\infty$ and all others one.
The contribution $-br+\log r$ diverges even when $b=0$.
Substituting Fenchel equality consequently gives the Bregman duality
$D_{\Phi_n}(X,Y)=D_{\Phi_n^*}(\Theta_Y,\Theta_X)$.

Finally $\Phi_n(X)\geq0$ gives the first determinant bound.  Apply
it to $nX/\operatorname{tr}X$ for the second and its equality case.
The Hessian of $\log\det$ is $-g$, which proves strict concavity on
the convex cone.
\end{proof}

The gradient in \eqref{final:matrix-gradient}, together with
$F(E_n)=0$, identifies a differentiable $F$ on the cone: the gradient
of $F-\Phi_n$ vanishes, so integration on each straight segment makes
the difference constant.  Likewise the full Hessian
\eqref{final:matrix-hessian}, together with $F(E_n)=0$ and
$\nabla F(E_n)=0$, identifies a twice differentiable $F$, since the
Hessian of the difference vanishes and its gradient is constant on
segments.  These affine marks remove exactly the affine ambiguity.
The rank-one potential or its calibrated anchored Bregman function
recovers $I$ directly.  A metric given merely up to isometry does not
provide these affine coordinates or marks.  Nor does congruence
invariance alone select this metric: the forms
\[
 \alpha g_X(U,V)+\beta\operatorname{tr}(X^{-1}U)
                             \operatorname{tr}(X^{-1}V),
 \quad \alpha>0,\quad\beta>-\alpha/n,
\]
are also invariant positive metrics.  Positivity follows by separating
$X^{-1/2}UX^{-1/2}$ into its scalar and traceless parts.  The selected
Hessian, rather than bare invariance or determinant inequalities,
removes this freedom.

\begin{theorem}[Geodesics, distance, and the symmetrization boundary]
\label{final:M3}
For $X,Y\in\mathcal P_n$ put
$K=\log(X^{-1/2}YX^{-1/2})$.  The unique constant-speed minimizing
geodesic for $g$ is
\[
 \gamma(s)=X^{1/2}e^{sK}X^{1/2},\quad0\leq s\leq1,
 \qquad d(X,Y)=\|K\|_F.
\]
If $\lambda_i$ are the eigenvalues of $X^{-1/2}YX^{-1/2}$, then
\[
 D(X,Y)+D(Y,X)=4\sum_i\sinh^2\left(\frac{\log\lambda_i}{2}\right).
\]
For $n=1$ this is $4\sinh^2(d(X,Y)/2)$.  For every $n\geq2$,
symmetrization is not a function of distance alone.
\end{theorem}

\begin{proof}
Congruence reduces the distance problem to a curve from $E_n$ to
$A=X^{-1/2}YX^{-1/2}$.  Write any piecewise continuously differentiable
such curve as $B(s)=e^{Z(s)}$ with $Z(s)$ symmetric.  The exponential
differential, obtained by differentiating its power series, is
\[
 D\exp_Z[V]=\int_0^1e^{(1-r)Z}Ve^{rZ}\,dr.
\]
In a basis diagonalizing $Z$ with eigenvalues $z_i$, its $(i,j)$ entry
is $V_{ij}(e^{z_i}-e^{z_j})/(z_i-z_j)$, with the continuous value
$e^{z_i}V_{ij}$ when $z_i=z_j$.  Consequently
\[
 g_B(B',B')=
 \sum_{i,j}\left(\frac{2\sinh((z_i-z_j)/2)}{z_i-z_j}\right)^2
                        (Z'_{ij})^2\geq\|Z'\|_F^2,
\]
where the quotient is one at zero.  Its absolute value is at least
one by the power series for $\sinh$.  Integrating speed gives
\[
 \operatorname{length}_g(B)\geq\int_0^1\|Z'(s)\|_F\,ds
                  \geq\|\log A\|_F.
\]
The curve $e^{s\log A}$ attains both inequalities and has constant
speed.  For a minimizing curve equality in the Euclidean length
inequality forces $Z(s)=a(s)\log A$ with $a$ nondecreasing from zero
to one.  On this curve metric speed is $a'(s)\|\log A\|_F$;
constant speed forces $a(s)=s$.  If $A=E_n$, zero length forces the
constant curve.  This proves uniqueness as stated.  The minimizing
curve is a geodesic by the usual first-variation theorem; directly
it satisfies $\gamma''=\gamma'\gamma^{-1}\gamma'$, the Euler--Lagrange
equation obtained from the energy
$\frac12\int\operatorname{tr}(\gamma^{-1}\gamma'
\gamma^{-1}\gamma')\,ds$.

Adding the two divergence formulas cancels logarithms and gives
$\sum_i(\lambda_i+\lambda_i^{-1}-2)$, proving symmetrization.
For $r>0$ compare $X=E_n$ with matrices whose log eigenvalue vectors
are $(r,0,\ldots,0)$ and
$(r/\sqrt2,r/\sqrt2,0,\ldots,0)$.  Both distances equal $r$, whereas
their symmetrizations are $2\cosh r-2$ and
$4\cosh(r/\sqrt2)-4$.  The difference has power series
\[
 \sum_{j\geq2}\frac{2-2^{2-j}}{(2j)!}r^{2j}>0.
\]
This proves the failure for every positive distance in rank at least two.
\end{proof}

\begin{theorem}[The supplied Gaussian and Wishart sampling model]
\label{final:M4}
Let $m\geq1$ and $n\geq1$ be specified integers and suppose
$Z_1,\ldots,Z_m$ are independent, identically distributed
$N(0,X)$ vectors with $X\in\mathcal P_n$ and known mean zero.  Put
$W=\sum_{j=1}^mZ_jZ_j^{\mathsf T}$ and $C=W/m$.  The scatter $W$ has
the Wishart law defined by this model and is a sufficient statistic
for $X$, with
\[
 \log L(X)=-\frac m2\log\det X
              -\frac12\operatorname{tr}(X^{-1}W)+\text{constant}.
\]
For symmetric $\Theta$,
\[
 \mathbb E e^{\operatorname{tr}(\Theta W)}
   =\det(E_n-2X^{1/2}\Theta X^{1/2})^{-m/2}
\]
when the matrix inside the determinant is positive definite; the
expectation is infinite otherwise.  The scatter is positive definite
almost surely exactly when $m\geq n$.  Whenever the observed $C$ is
positive definite, the maximum likelihood covariance is uniquely $C$
and, with $\ell=-\log L$,
\[
 \ell(X)-\ell(C)=\frac m2D(C,X).
\]
If $C$ is singular there is no maximum likelihood estimate in the
open cone $\mathcal P_n$.  Finally
\[
 \operatorname{KL}(N(0,X)\,\|\,N(0,Y))=\frac12D(X,Y),
\]
and the divergence for $m$ independent copies is $m$ times this value.
\end{theorem}

\begin{proof}
Multiplying the Gaussian densities proves the likelihood formula.
Its dependence on the samples factors through $W$, so the
factorization criterion for dominated statistical families proves
sufficiency.  For one sample write $Z=X^{1/2}G$ with $G\sim N(0,E_n)$.
The tilted Gaussian integral has quadratic form
$E_n-2X^{1/2}\Theta X^{1/2}$.  Orthogonal diagonalization reduces it
to one-dimensional integrals, giving its inverse square-root
determinant when all eigenvalues are positive.  A zero or negative
eigenvalue yields an infinite one-dimensional integral, so Tonelli's
theorem gives divergence otherwise.  Independence raises the finite
answer to the $m$th power.

The rank of $W$ is the dimension of the span of its $m$ sample vectors.
It is at most $m$, and the first $n$ Gaussian sample vectors are linearly
independent almost surely: the determinant-zero set has Lebesgue measure
zero, as follows, for example, by conditioning successively on the span
of the preceding vectors. Thus $W$ is positive definite almost surely
for $m\geq n$, and singular for $m<n$.

For positive-definite $C$, subtracting likelihood values gives
$mD(C,X)/2$, with the asserted minimum and equality case. For singular
$C$, diagonalize it and let the eigenvalues of $X$ along the nontrivial
kernel of $C$ tend to zero while keeping the remaining eigenvalues fixed
and positive. The trace term stays fixed and $\log\det X\to-\infty$,
so $\ell(X)\to-\infty$ and no minimizer exists in the open cone.
Finally integrating the Gaussian log
density ratio uses $\mathbb E_X ZZ^{\mathsf T}=X$ and gives $D(X,Y)/2$;
the log ratio adds over independent samples.
\end{proof}

The sample count, independence, Gaussian law, and known-mean convention
are data of this theorem.  The scalar potential alone supplies none
of those sampling choices.  Likewise the spectral or Gaussian
construction of $\Phi_n$ identifies that constructed family; an
independently designated family is identified only after the rules
of Theorem~\ref{final:M1} are imposed.

\section{Gaussian radial realization and independent spatial data}
\label{final:spatial}

\begin{theorem}[The isotropic four-dimensional Gaussian realization]
\label{final:R1}
If $Z\sim N(0,E_4)$ and $T=\|Z\|^2/2$, then $T$ has density
$p(t)=te^{-t}$ on $t>0$.  Conversely, if $T$ has this density and
$\Omega$ is an independent uniform point of $S^3$, then
$Z=\sqrt{2T}\,\Omega$ has law $N(0,E_4)$.  Equivalently the radial
law together with orthogonal invariance uniquely identifies this
Gaussian law; the radial law alone does not.

For the specified Ornstein--Uhlenbeck process in $\mathbb R^D$,
\[
 dZ_s=-\tfrac12Z_s\,ds+dB_s,
\]
the generator on radial-energy functions $F(z)=f(\|z\|^2/2)$ is
\[
 \left(\tfrac12\Delta-\tfrac12z\cdot\nabla\right)F
       =t f''(t)+(D/2-t)f'(t).
\]
Thus its radial generator is $tf''+(2-t)f'$ exactly when $D=4$.
In dimension four the radial energy satisfies
\[
 dT_s=(2-T_s)\,ds+\sqrt{2T_s}\,d\beta_s
\]
for a one-dimensional Brownian motion $\beta$.  These statements
select a dimension within this specified Gaussian construction.
\end{theorem}

\begin{proof}
Polar integration in $\mathbb R^4$, using $|S^3|=2\pi^2$, gives
\[
 (2\pi)^{-2}e^{-r^2/2}\,2\pi^2r^3\,dr
       =te^{-t}\,dt,\qquad t=r^2/2.
\]
This proves the forward law.  Conversely, integration of any bounded
Borel test function first over uniform $\Omega$ and then over this
radial density reproduces the four-dimensional Gaussian integral.
For an orthogonally invariant law, average the test function over
the normalized Haar measure on the compact group $O(4)$.  On each
nonzero sphere this average is its uniform spherical average, so
its integral is fixed by the radial law.  This proves the equivalent
uniqueness statement and the independence of radius and uniform
angle.  Choosing instead a fixed unit vector $v$ and
$Z=\sqrt{2T}\,v$ preserves the radial law and changes the vector law.

The chain rule gives
$\nabla F=f'(t)z$, $\Delta F=2tf''(t)+Df'(t)$, and
$z\cdot\nabla F=2tf'(t)$, proving the generator formula.
It\^o's formula yields
\[
 dT_s=(D/2-T_s)\,ds+Z_s\cdot dB_s,
 \qquad d\langle T\rangle_s=2T_s\,ds.
\]
Normalize the stochastic integrand to $Z_s/\|Z_s\|$ off zero and
choose a fixed unit vector at zero.  Its stochastic integral
$\beta_s$ has quadratic variation $s$, hence is Brownian by the
continuous-martingale characterization, and
$Z_s\cdot dB_s=\sqrt{2T_s}\,d\beta_s$.  At $D=4$ this gives the
displayed equation.  The stationary vector law is $N(0,E_D)$ by the
explicit Gaussian solution of the OU equation; polar integration
therefore gives the stationary energy law
$t^{D/2-1}e^{-t}/\Gamma(D/2)$.  Its shape is two precisely at $D=4$.
\end{proof}

This construction retains the Gaussian and stochastic-calculus
foundations just used.  A stationary density by itself does not
choose a diffusion: for any specified positive differentiable
coefficient $a(t)$ the formal zero-current expression
\[
 L_af=\frac1p(apf')'
     =af''+\left(a'+a\frac{p'}p\right)f'
\]
has the same formal symmetric measure, as integration by parts on
compactly supported tests verifies.  Domain and endpoint conditions
remain necessary when selecting an actual closed generator.  Choosing
$a(t)=t$ gives the radial operator in Theorem~\ref{final:R1}; it is a
selection, not a consequence of knowing only $p$.

\begin{theorem}[Nonnegative orthogonal additivity without regularity]
\label{final:R2}
Let $V$ be a real inner-product space of dimension at least two and
let $r:V\to[0,\infty)$.  Then
\[
 r(x+y)=r(x)+r(y)\quad\text{whenever }\langle x,y\rangle=0
\]
holds if and only if $r(x)=c\|x\|^2$ for some $c\geq0$.  No continuity,
measurability, homogeneity, or radiality hypothesis is required.
Consequently a function $z:V\to[0,\infty)$ satisfies
\[
 z(x+y)=z(x)+z(y)+z(x)z(y)\quad(x\perp y)
\]
if and only if $z(x)=e^{c\|x\|^2}-1$ for some $c\geq0$.
Nontriviality forces $c>0$, while a numerical value at one nonzero
vector is needed to calibrate its scale.
\end{theorem}

\begin{proof}
The equation at $(0,0)$ gives $r(0)=0$.  Define the even and odd parts
\[
 e(x)=\frac{r(x)+r(-x)}2,\qquad
 o(x)=\frac{r(x)-r(-x)}2.
\]
They are orthogonally additive, $e\geq0$, and $|o(x)|\leq e(x)$.
If $\|x\|=\|y\|$, set $a=(x+y)/2$, $b=(x-y)/2$.  Then $a\perp b$,
$x=a+b$, and $y=a-b$.  Evenness gives
$e(x)=e(a)+e(b)=e(a)+e(-b)=e(y)$.  Every nonnegative real number is
the squared norm of a vector, so $e(x)=A(\|x\|^2)$ for a function
$A:[0,\infty)\to[0,\infty)$.

There exist two orthogonal unit vectors.  Multiplying them by
$\sqrt a$ and $\sqrt b$ shows $A(a+b)=A(a)+A(b)$ for arbitrary
$a,b\geq0$.  Nonnegativity makes $A$ monotone: if $b\geq a$ then
$A(b)-A(a)=A(b-a)\geq0$.  With $c=A(1)$, additivity gives
$A(q)=cq$ for every nonnegative rational $q$.  For any real $a\geq0$,
choose rational sequences approaching $a$ from below and above.
Monotonicity sandwiches $A(a)$ between their multiples by $c$, so
$A(a)=ca$.

For $x\ne0$, choose $y\perp x$ of the same norm; the case $x=0$ is
automatic.  Then $x+y\perp x-y$, and orthogonal additivity and oddness
give
\[
 o(2x)=o(x+y)+o(x-y)
      =o(x)+o(y)+o(x)+o(-y)=2o(x).
\]
Iterating downward, $o(x)=2^ko(x/2^k)$, while
$|o(x/2^k)|\leq c\|x\|^2/4^k$.  Hence
$|o(x)|\leq c\|x\|^2/2^k$ for every $k$, forcing $o(x)=0$.
This proves $r=e=c\|x\|^2$.  Pythagoras proves the converse.

For the second statement take $r=\log(1+z)$.  Its values are
nonnegative, and the composition equation is exactly orthogonal
additivity of $r$.  Apply the first statement and exponentiate;
the converse follows by multiplication of exponentials.
\end{proof}

Both substantive hypotheses have direct omission tests.  In dimension
one, $r(x)=x^4$ is nonnegative and orthogonally additive, since every
orthogonal pair includes zero, but is not $c\|x\|^2$.  Without
nonnegativity in dimension at least two, a nonzero linear functional,
or $\|x\|^2$ plus such a functional, is orthogonally additive and
violates the conclusion.  Accordingly the condition $z>-1$ alone
also fails: $z(x)=e^{\ell(x)}-1$ for a nonzero linear functional
obeys the composition rule and takes negative values.

\begin{proposition}[The profile-independent radial residual]
\label{final:R3}
In integer Euclidean dimension $D\geq1$ the radial Laplacian is
$\Delta_D f=f''+(D-1)f'/x$ for $x>0$.  For every twice differentiable
nonvanishing profile $f$ and $\alpha=(D-1)/2$,
\begin{equation}\label{final:radial-residual}
 \frac{(x^\alpha f)''}{x^\alpha f}
    -\left(\frac{f''}{f}+\frac{D-1}{x}\frac{f'}{f}\right)
     =\frac{(D-1)(D-3)}{4x^2}.
\end{equation}
Equivalently, writing $f=x^{-\alpha}u$ gives the operator identity
\[
 -\Delta_D f=x^{-\alpha}
   \left(-u''+\frac{(D-1)(D-3)}{4x^2}u\right),
\]
which does not require $f$ to be nonvanishing.  Residual cancellation
holds exactly for $D=1$ or $D=3$; with the independently retained
restriction $D\geq2$ it selects $D=3$.  It has no reverse identifying
force for the profile $f$.
\end{proposition}

\begin{proof}
Twice applying the product rule gives
\[
 \frac{(x^\alpha f)''}{x^\alpha f}
 =\frac{f''}{f}+\frac{2\alpha}{x}\frac{f'}{f}
                           +\frac{\alpha(\alpha-1)}{x^2}.
\]
Since $2\alpha=D-1$ and
$\alpha(\alpha-1)=(D-1)(D-3)/4$, subtraction proves
\eqref{final:radial-residual}; the same product rule proves the
operator form.  The coefficient vanishes precisely at the two stated
integer dimensions.  It is independent of $f$, so for example
$f(x)=e^{-x^2}$ also gives cancellation in dimension three, despite
being a different profile from $p$ or the scalar potential.
\end{proof}

\begin{proposition}[Spatial countermodels retaining the complete intrinsic structure]
\label{final:R4}
All intrinsic scalar identities, their probability and operator
realizations, their formal characteristic series, and the canonical
four-dimensional Gaussian construction may hold while an independent
designated spatial observable fails orthogonal additivity, or while
its designated dimension fails to equal three.  Their conjunction therefore
does not imply either the hypotheses of Theorem~\ref{final:R2} or a
designated dimension conclusion from Proposition~\ref{final:R3}.
\end{proposition}

\begin{proof}
Keep the entire intrinsic structure and every canonical auxiliary
object fixed.  Adjoin a distinct Euclidean spatial sort $V$ with its
own designated observable $r$.  Three choices give the required
models.  In $V=\mathbb R^2$ take $r(x)=\|x\|^4$.  For perpendicular
unit vectors $x,y$, $r(x+y)=4$ while $r(x)+r(y)=2$, so even positivity,
smoothness, and radiality do not force the composition law.  In
$V=\mathbb R^2$ take $r(x)=\|x\|^2$.  Orthogonal additivity holds,
but the coefficient in \eqref{final:radial-residual} is $-1/(4x^2)$.
Finally in $V=\mathbb R$ take squared norm: residual cancellation
holds although $D\ne3$.  None of these additions changes any object
or equation in the original intrinsic structure.  In particular,
each coexists with the separately constructed Gaussian vector in
$\mathbb R^4$; no identification between its vector space and the
designated spatial sort was supplied.
\end{proof}

There is also a direct tensor distinction.  Pulling a one-dimensional
metric $g(r)\,dr^2$ back by a spatial scalar map produces
$g(r)\,dr\otimes dr$, of rank at most one.  The Hessian of a composite
potential is instead
\[
 D^2(I\circ r)=I''(r)\,dr\otimes dr+I'(r)D^2r.
\]
The chain rule proves this formula.  Its second term prevents an
identification of these two tensors or an inference of ambient
dimension from the rank of the first.

\section{Marked rational dynamics and transported integer arithmetic}
\label{final:arithmetic}

\begin{theorem}[Rational trees, Euclidean decoding, and the matrix monoid]
\label{final:B1}
Fix the positive rational coordinate, the root $1$, ordered edge
labels $L,R$, and the transformations
\[
 L(x)=\frac{x}{1+x},\quad R(x)=1+x,\qquad
 M_L=\begin{pmatrix}1&0\\1&1\end{pmatrix},\quad
 M_R=\begin{pmatrix}1&1\\0&1\end{pmatrix}.
\]
Successively applying the letters of a word to $1$ gives a rooted
Calkin--Wilf tree~\cite{CalkinWilf} in which every positive rational occurs exactly
once.  The matrices $M_L,M_R$ freely generate exactly the determinant-one
matrices with nonnegative integer entries, including the identity.

For a Stern--Brocot interval~\cite{Stern1858} store the upper endpoint as the first
column and the lower endpoint as the second, starting with columns
$(1,0)^{\mathsf T}$ and $(0,1)^{\mathsf T}$, representing $\infty$ and
$0$.  Its node is their mediant and a left or right step multiplies
this matrix on the right by $M_L$ or $M_R$.  With these conventions
the Stern--Brocot label of a word is the Calkin--Wilf label of the
reversed word.

The complete root path, its oriented subtraction runs, and the
finite continued fraction are mutually decodable, with the canonical
convention
\[
 [a_0;a_1,\ldots,a_k],\quad a_0\geq0,\quad
 a_j\geq1\ (j\geq1),\quad a_k\geq2\ \text{if }k\geq1.
\]
A one-term expansion $[a_0]$ has $a_0\geq1$, so the root is $[1]$.
The last run in the subtraction-to-root algorithm has length one less
than its terminal Euclidean quotient; it has length zero for the root.
\end{theorem}

\begin{proof}
For a reduced fraction $a/b<1$, the unique possible parent is
$a/(b-a)$ and the last edge is $L$; if $a/b>1$, it is $(a-b)/b$ with
last edge $R$.  Both numerators and denominators remain positive,
the gcd is unchanged, and their sum strictly decreases.  The
algorithm terminates only at $a=b$, hence at $1/1$ by coprimality.
It therefore constructs a root path for every positive rational.
At each step both the last edge and parent are forced, proving
uniqueness.  Conversely the child formulas preserve positive
coprimality, so every word has a positive reduced rational label.

For matrix generation write a nonnegative determinant-one matrix as
$M=\left(\begin{smallmatrix}a&b\\c&d\end{smallmatrix}\right)$.
Except for the identity, its two rows are comparable coordinatewise.
Indeed one pattern of incomparable rows, $a<c$ and $b>d$, would
give $ad<bc$.  In the other pattern write $a=c+u$, $d=b+v$ with
$u,v\geq1$.  The determinant equation becomes
\[
 1=cv+ub+uv,
\]
forcing $u=v=1$ and $b=c=0$, which is the identity.  If the lower
row dominates, subtract the upper row to obtain $M=M_LN$; if the
upper dominates, subtract the lower to obtain $M=M_RN$.
The resulting $N$ remains nonnegative integral with determinant one
and strictly smaller entry sum.  Equal rows are impossible.  Repeated
subtraction therefore ends at the identity and proves generation.
The dominance choice at every nonidentity step is unique; cancellation
then proves uniqueness of the generator word.  Conversely products
have nonnegative integral entries and determinant one.  The word
``nonnegative'' is essential: the generators have zero entries.

If $w=s_1\cdots s_\ell$ records chronological path steps, its
Calkin--Wilf column vector is
$M_{s_\ell}\cdots M_{s_1}(1,1)^{\mathsf T}$.  The Stern--Brocot
endpoint matrix is $M_{s_1}\cdots M_{s_\ell}$ and its mediant vector
is that matrix times $(1,1)^{\mathsf T}$.  This proves reversal and
also proves complete Stern--Brocot enumeration.  For example,
$\operatorname{CW}(LR)=3/2$ and $\operatorname{SB}(LR)=2/3$.
Right multiplication by $M_L$ replaces the upper endpoint column
by its sum with the lower, and right multiplication by $M_R$ replaces
the lower by that sum.  Determinant one ensures ordered adjacent
endpoints and a reduced mediant throughout.

Lastly, consecutive subtractions of the same smaller integer from
the larger group into Euclidean integer division.  If $a=qb+r$ with
$0<r<b$, exactly $q$ subtractions replace $(a,b)$ by $(r,b)$.
If $r=0$, coprimality makes $b=1$ and the algorithm stops after
$q-1$ subtractions at $(1,1)$; the case with $b>a$ is identical
with the roles reversed.  Thus oriented runs determine every
quotient, with precisely the stated terminal correction.  The
Euclidean algorithm terminates and its remainders are unique.
Successive identities $a/b=q+r/b$ yield the continued fraction.
If its last digit is one, the identity
$[\ldots,c,1]=[\ldots,c+1]$ removes it; the unique division
quotients give the stated canonical terminal convention.  A rational
value reconstructed from that fraction gives its unique root path
by the already proved parent algorithm, which proves mutual decoding.
\end{proof}

\begin{proposition}[Marked Farey branches and full-domain scalar recovery]
\label{final:B2}
On $[0,1]$ define the Farey map and its inverse branches~\cite{IsolaFarey}
\[
 \psi_0(x)=\frac{x}{1+x},\qquad \psi_1(x)=\frac1{1+x}.
\]
Their respective ranges are $[0,1/2]$ and $[1/2,1]$, and they are
the increasing and decreasing inverse branches of
\[
 F(y)=
 \begin{cases}
 y/(1-y),&0\leq y\leq1/2,\\
 (1-y)/y,&1/2\leq y\leq1.
 \end{cases}
\]
At the common endpoint both forward formulas equal one, and branch
labels retain the inverse-branch information.

The complete numerically labelled Calkin--Wilf tree determines $L,R$
on positive rationals.  If the candidate maps are declared to be
M\"obius transformations, these data determine their full rational
formulas, in particular $L$ on $(-1,\infty)$.  In that marked
extension category the scalar potential is recovered by
\begin{equation}\label{final:tree-integral}
 I(t)=\int_0^{t-1}L(s)\,ds=t-1-\log t,\qquad t>0.
\end{equation}
Connected real analytic continuation is another sufficient extension
category.  A bare tree, or positive-rational action without a suitable
extension category, has no such global converse.
\end{proposition}

\begin{proof}
Derivatives, endpoint values, and substitution verify the ranges,
monotonicity, and the two inverse formulas.  Notice that $\psi_1$ is
not the Calkin--Wilf map $R$.  Every positive rational occurs in the
tree, so its numerical children give both maps on all such inputs.
Two M\"obius transformations agreeing at three distinct finite inputs
where they are defined are equal: cross multiplication gives a
polynomial of degree at most two with three roots, hence the zero
polynomial.  Their fractional-linear maps therefore agree globally
away from their pole.  This supplies the extension $s/(1+s)$ for
$-1<s<0$ as well as for positive $s$, and direct integration proves
\eqref{final:tree-integral}.  Within the real analytic category on
$(-1,\infty)$, equality on positive rationals first gives equality
on $(0,\infty)$ by continuity and then everywhere by the identity
theorem.

Without that category, choose a nonzero smooth bump $\chi$ supported
in $(1/4,1/2)$ and put $\widetilde I=I+\varepsilon\chi$.
The induced generator $\widetilde L(s)=\widetilde I'(1+s)$ agrees
with $L(s)$ for every $s\geq0$, and hence on all positive-rational
inputs, but the potential differs on $(0,1)$.  Small nonzero
$\varepsilon$ even preserves strict convexity, as in
Proposition~\ref{final:F3}.  Finally an abstract unlabelled binary
tree contains no numerical input/output pairs at all, so it cannot
supply the coordinate markings used in this proof.
\end{proof}

The scalar formal-group inverse itself constructs the canonical
global rational expression: $\operatorname{inv}_\star(x)=-x/(1+x)$,
so $L=-\operatorname{inv}_\star$ and $R(x)=1+x$ is the marked shift.
This constructs a labelled rational model.  It does not attach those
labels to an independently supplied abstract tree.

\begin{theorem}[Exact collisions and complete labelled arithmetic transport]
\label{final:B3}
Fix $\mu,\gamma>0$ and a real constant $c_P$, and consider the placed
integer code
\[
 e(n)=\log(1+\gamma n)-\mu n+c_P,\qquad n\geq2.
\]
For distinct $n>m\geq2$,
\begin{equation}\label{final:integer-collision}
 e(n)=e(m)\quad\Longleftrightarrow\quad
 \mu=\frac{1}{n-m}\log\frac{1+\gamma n}{1+\gamma m}.
\end{equation}
The rule $e(a)\odot e(b)=e(ab)$ defines a well-defined binary
operation on the image if and only if $e$ is injective.  Under that
hypothesis, write $S$ for the image with an adjoined abstract unit $1_S$.
Extending $e(1)=1_S$ makes $e:\mathbb N_{>0}\to S$ an isomorphism of labelled multiplicative
monoids.  It transports the full prime, valuation, divisor,
gcd/lcm, M\"obius, and Dirichlet-convolution inventory stated below.

The different intrinsic code
\[
 e_0(n)=H(n+1),\qquad n\geq1,
\]
is injective for every placement and already includes the unit label.
It does not remove any collision of the original code $e$.
\end{theorem}

\begin{proof}
Subtracting the two values of $e$ proves
\eqref{final:integer-collision}; thus avoiding every displayed
pairwise parameter curve is exactly injectivity on the specified
carrier.  If injective, $e$ is a bijection onto its image and ordinary
multiplication transports through that bijection, proving
well-definedness.

For the reverse implication suppose $n>m$ and $e(n)=e(m)$.  Then
\[
 \mu(n-m)=\int_m^n\frac{\gamma}{1+\gamma u}\,du.
\]
For every real $k>1$, the same difference calculation and substitution
in the integral give
\begin{align*}
 e(kn)-e(km)
  &=k\int_m^n\left(\frac{\gamma}{1+k\gamma u}
                        -\frac{\gamma}{1+\gamma u}\right)du<0.
\end{align*}
Here the formula for $e$ is evaluated on the positive real ray in
this calculation.  The strict inequality follows pointwise from
$u>0$, $\gamma>0$, and $k>1$.  Taking an integer $k\geq2$ shows that
multiplication of the common image $e(n)=e(m)$ by $e(k)$ would
require two distinct outputs.  Hence no operation with the prescribed
rule exists.  This proves necessity for every collision, not only
for a selected example.  For instance, $\gamma=1$ and
$\mu=\log(5/4)$ give $e(3)=e(4)$ but
$e(6)-e(8)=\log(175/144)>0$.

When injectivity holds, adjoin an element $1_S$ distinct from every
point of the original image and set $1_S\odot x=x\odot1_S=x$.
The extended map is a bijection preserving multiplication and unit.
The integer fundamental theorem of arithmetic consequently gives a
unique finite factorization of each $x\in S$ into irreducibles,
with the empty product for the unit.  All the inventory formulas
below follow from that factorization, as verified after the table.
Finally $H'(t)=1/t-1<0$ for $t>1$, so $H(n+1)$ is strictly decreasing
on positive integers.  This proves the alternative's injectivity.
For the placed profile with arbitrary additive level, write
$\widetilde\Sigma_P(r)=\log(1+\gamma r)-\mu r+c_P$;
its parameters here are $\mu,\gamma>0$ and $c_P\in\R$.
The critical point is
$r_*=1/\mu-1/\gamma$ and
\[
 \widetilde\Sigma_P(r_*+n/\mu)=H(n+1)
                   +\log(\gamma/\mu)+\mu/\gamma-1+c_P.
\]
Indeed its intrinsic coordinate $\mu(r+1/\gamma)$ is $n+1$ there.
This is a shifted sampling code, with a fixed additive level, rather
than the original sampling at $r=n$.
\end{proof}

Write $x\mid y$ for divisibility in the unit-adjoined monoid, let
$q$ range over its irreducibles, and write $\mu_S$ for the arithmetic
M\"obius function, to distinguish it from the placement parameter.
The complete inventory is
\[
\begin{array}{ll}
 \text{Irreducibles}&q=e(p),\quad p\text{ a numerical prime};\\[2pt]
 \text{Valuations}&v_q(x)=\max\{j\geq0:q^j\mid x\};\\[2pt]
 \text{Factor counts}&\Omega(x)=\sum_qv_q(x),\quad
             \omega(x)=\#\{q:v_q(x)>0\};\\[2pt]
 \text{Greatest common divisor}&
       v_q(\gcd(x,y))=\min(v_q(x),v_q(y));\\[2pt]
 \text{Least common multiple}&
       v_q(\operatorname{lcm}(x,y))=\max(v_q(x),v_q(y));\\[2pt]
 \text{Divisor lattice}&
       \{d:d\mid x\}\cong\prod_{q\mid x}\{0,\ldots,v_q(x)\};\\[2pt]
 \text{Divisor count}&\tau(x)=\prod_{q\mid x}(v_q(x)+1);\\[2pt]
 \text{M\"obius function}&
       \mu_S(x)=0\ \text{if some }v_q(x)\geq2,
       \quad\mu_S(x)=(-1)^{\omega(x)}\ \text{otherwise};\\[2pt]
 \text{Dirichlet convolution}&
       (f*g)(x)=\sum_{d\mid x}f(d)g(x/d).
\end{array}
\]
All sums and products over the support of a single element are finite,
and $\mu_S(1_S)=1$.  To verify the inventory, divisibility of
$x=\prod_q q^{v_q(x)}$ by $d$ is precisely the componentwise condition
$0\leq v_q(d)\leq v_q(x)$.  This proves the valuation maximum,
gcd/lcm formulas, lattice product, and divisor count.  Irreducibility
means exactly that the exponent vector has one coordinate equal to
one and all others zero; it therefore identifies the numerical
primes under $e$.  Summing $\mu_S$ over divisors leaves only the
square-free choices, giving
\[
 \sum_{d\mid x}\mu_S(d)=(1-1)^{\omega(x)}
 =\begin{cases}1,&x=1_S,\\0,&x\ne1_S.\end{cases}
\]
In any fixed commutative coefficient ring, finite reindexing of
divisor pairs makes convolution associative and commutative with
identity $\delta_{1_S}$.  The last displayed identity says
$\mu_S*\mathbf1=\delta_{1_S}$.  Thus if
$F(x)=\sum_{d\mid x}f(d)$, convolution with $\mu_S$ gives
\[
 f(x)=\sum_{d\mid x}\mu_S(d)F(x/d),
\]
proving M\"obius inversion.  A normalized multiplicative function
satisfies $f(1_S)=1$ and $f(xy)=f(x)f(y)$ for coprime $x,y$.
It is exactly specified by arbitrary values at the marked prime
powers, followed by
$f(x)=\prod_{q\mid x}f(q^{v_q(x)})$.  The operation does not choose
those values.  Over a field, every nonzero function satisfying the
unnormalized multiplicativity equation automatically has unit
value one; the identically zero function is a separate convention.

On the original image omitting one, reflexive divisibility means
$x=y$ or $y=x\odot z$ with $z$ in that image.  Coprime gcds and the
bottom of a full divisor lattice require the adjoined $1_S$.  It
cannot be identified with the real number $\widetilde\Sigma_P(1)$ unless
injectivity has also been checked after adding that index.
Prime permutations preserve the abstract free commutative monoid,
so numerical prime names and addition are extra labels.  With the
integer labels supplied, the totient is also transported by
\[
 \varphi(e(n))=\prod_{p^a\parallel n}p^{a-1}(p-1),
 \qquad\varphi(1_S)=1.
\]
The formula counts residues not divisible by any prime divisor by
inclusion--exclusion: for each square-free divisor $d$ of $n$ there
are $n/d$ multiples of $d$ among $1,\ldots,n$, giving
$n\prod_{p\mid n}(1-1/p)$.  Its numerical factors show why this
function requires the integer labels and their sizes.  No step in
this arithmetic transport identifies a real embedding or the
analytic scalar shape from an unlabelled monoid.

\begin{theorem}[Numerical Euler products and the full multiset inverse]
\label{final:B4}
With the numerical positive-integer labels and their usual
multiplication,
\[
 \zeta(s)=\sum_{n\geq1}n^{-s}
          =\prod_{p\text{ prime}}(1-p^{-s})^{-1},
 \qquad\operatorname{Re}s>1.
\]
If a supplied nonnegative self-adjoint operator $A$ has a complete
simple eigenbasis with eigenvalues $0,1,2,\ldots$, then on that same
domain $\operatorname{Tr}(1+A)^{-s}=\zeta(s)$.

Conversely, let a multiset of real generators $q>1$ have positive
integer multiplicities $m_q$.  Suppose its Euler product
\[
 Z(s)=\prod_q(1-q^{-s})^{-m_q}
\]
converges to a finite positive value for all sufficiently large real
$s$.  The full real function on any such terminal ray uniquely
determines the numerical multiset, including multiplicities.
For a nonempty multiset its least generator $q_0$ and multiplicity
$m_0$ are recovered by
\[
 q_0=\lim_{s\to\infty}(\log Z(s))^{-1/s},\qquad
 m_0=\lim_{s\to\infty}q_0^s\log Z(s).
\]
Removing the entire factor $(1-q_0^{-s})^{-m_0}$ and repeating
recovers every remaining generator.  The empty multiset is exactly
$Z\equiv1$.
\end{theorem}

\begin{proof}
Expand the geometric factors over any finite set of primes.  Unique
factorization gives the sum over exactly the integers supported on
those primes.  Absolute convergence of
$\sum_{n\geq1}n^{-\operatorname{Re}s}$ permits the limit over all
primes and proves the identity.  In the indicated eigenbasis the
operator $(1+A)^{-s}$ is diagonal with entries $(n+1)^{-s}$ and
singular values $(n+1)^{-\operatorname{Re}s}$.  It is trace class
precisely for $\operatorname{Re}s>1$, with the asserted trace.

For the inverse choose a sufficiently large $s_0>0$ at which the
product converges.  Positivity of its factors gives
\[
 \log Z(s_0)=\sum_qm_q[-\log(1-q^{-s_0})]
           =\sum_qm_q\sum_{k\geq1}\frac{q^{-ks_0}}k<\infty.
\]
Tonelli's theorem justifies the expansion.  For each finite $M>1$
every generator $q\leq M$ contributes at least $M^{-s_0}$, counted
with multiplicity.  There are therefore only finitely many such
generators.  In particular the multiset is locally finite down to
one, is countable, and if nonempty has a least $q_0>1$ of finite
multiplicity $m_0$.

For $s\geq s_0$ multiply the logarithmic series by $q_0^s$.
Every term is
$(m_q/k)(q_0/q^k)^s$ and is bounded by its value at $s_0$.
Those bounds sum to $q_0^{s_0}\log Z(s_0)<\infty$.  The ratio
$q_0/q^k$ equals one exactly when $q=q_0,k=1$, and is otherwise
strictly less than one.  Dominated convergence consequently gives
\[
 q_0^s\log Z(s)\longrightarrow m_0,
 \qquad\log Z(s)=m_0q_0^{-s}(1+o(1)).
\]
Taking the negative $1/s$ power gives the first limit and the second
is already proved.  Divide $Z(s)$ by the entire recovered factor,
or equivalently subtract
$-m_0\log(1-q_0^{-s})$ from its logarithm.  The remaining product
still satisfies the same convergence hypothesis.  This removes all
powers of $q_0$ at once, including those that coincide numerically
with another generator.  Repeating is therefore valid even in the
presence of such coincidences.  Every specified generator has only
finitely many predecessors by the bounded-interval argument, so
the iteration reaches all of them.  If none remains the logarithm
is zero; if one remains it is strictly positive.  This proves both
uniqueness and the empty-product assertion.
\end{proof}

The inverse in Theorem~\ref{final:B4} retains the positive Euler-factor
category and the numerical exponential coordinate $q^{-s}$.  It
does not reconstruct an addition law from abstract factors or
resolve any collision in Theorem~\ref{final:B3}.  Across these
sections, every construction and converse has its mathematical
proof and its retained foundations stated explicitly. The series,
topology, matrix, stochastic, and arithmetic development is included
in the complete Lean~4 formalization with no admitted proofs.
